% 本文件由 LaTeX 排版 Agent 根据有效 translation.json 自动生成。
% author=Councils; work_id=3813; title=君士坦丁堡第三届大公会议

\chapter{君士坦丁堡第三届大公会议}

% structure: p0001 source=h1 target=omitted reason=page-title-already-rendered-by-parent

% structure: p0002 source=h2 target=section reason=relative-heading-rank; base=h2

\section{第一届会议——会议纪要摘录}

\EditorialNote{\Emph{在叙述了会议召集的历史之后，会议纪要始于教宗代表的发言如下：}}\par

最仁慈的主上，遵照您受天主训诲的陛下致我们至圣教宗的诏书，我们受他派遣，来到您蒙天主确认的安详之至圣足前，携带他的\OriginalTerm{上呈文书}{\GreekText{ἀναφορᾶς}, suggestione}以及他的主教会议之另一上呈文书，后者同样由该会议属下可敬的主教们致您蒙天主护佑的虔诚。我们将这些文书呈献给了您蒙天主加冕的刚毅。然而，过去约四十六年间，某些与正统信仰相悖的表达新词，由曾是您这座皇城及蒙天主保守之城的主教们引入，即：\Person{Sergius}、\Person{Paul}、\Person{Pyrrhus}、\Person{Peter}，以及一度是亚历山大里亚城总主教的\Person{Cyrus}，还有曾是法蓝城主教的\Person{Theodore}，以及其某些追随者。这些事给普世教会带来了不小的混乱，因为他们教条式地教导：在我们主耶稣基督——天主圣三之一——降生成人的救世工程中，只有一个意志和一个操作。我们宗座——您的仆人——曾多次与之抗争，并为此祈祷，却至今未能使这种堕落意见的拥护者远离它。我们恳求您蒙天主加冕的刚毅，让那些与至圣君士坦丁堡教会持相同看法的人告诉我们，这种新奇说法的来源是什么。\par

\EditorialNote{\OriginalTerm{一志论者}{Monothelites}奉皇帝之命所作的回答：}\par

我们没有提出任何新的说法，而是教导了我们从神圣大公会议、从神圣公认的教父们、以及从这座帝都的总主教们——即\Person{Sergius}、\Person{Paul}、\Person{Pyrrhus}、\Person{Peter}——还有旧罗马教宗\Person{Honorius}和亚历山大里亚教宗\Person{Cyrus}那里所领受的一切，即关于意志和操作的教导。我们如此相信，如今仍如此相信，如此宣讲；并且我们准备坚守并捍卫这一信仰。\par

% structure: p0007 source=h2 target=section reason=relative-heading-rank; base=h2

\section{旧罗马教宗阿加托致皇帝的书信，以及阿加托与罗马主教会议一百二十五位主教致第六届大公会议的书信}

\EditorialNote{于十一月十五日第四届会议上应君士坦丁堡宗主教\Person{George}及其附属主教们的要求宣读。}\par

% structure: p0009 source=h3 target=subsection reason=relative-heading-rank; base=h2

\subsection{引言注释}

(Bossuet, \WorkTitle{Defensio Cler. Gal.} Lib. VII., cap. xxiv.)\par

所有教父逐一发言，只有在审查之后，圣阿加托的书信和整个西方主教会议才获得批准。的确，阿加托和西方主教们如此颁布他们的法令：\Quote{我们已从我们的卑微者中派遣了一些人，到你蒙天主保护的勇毅那里，他们将向你呈上我们所有人的报告，即所有北方或西方地区主教们的报告，其中我们也总结了我们宗徒信仰的宣认，但并非作为那些想要争论这些不确定之事的人，}而是作为确定不变之事，以简短的定义将其呈明，\Quote{谦卑地恳求你，借你神圣陛下的恩宠，命令将同一道理向所有人宣讲，并对所有人具有效力。}毫无疑问，就他们而言，他们已经定义了此事。问题是其他各教会是否同意；如此重大的事只有在主教审查后才得以澄清。但圣阿加托关于其宗座所用之崇高、宏伟而真实的表达，即：依靠主的许诺，它从未偏离真理之路，以及其教宗——阿加托的前任们——受命在伯多禄之位中坚固弟兄们，并始终履行这一职责，这些都被大公会议教父们听闻并接受。但他们同样审查此事，探究罗马教宗的法令，并在审查后批准阿加托的法令，谴责\Person{Honorius}的法令：这明确证明，他们并不将阿加托的表达理解为必须不经讨论就接受罗马教宗的每一项法令，即使是\OriginalTerm{信理}{de fide}方面的，因为这些法令受制于大公会议的最高最终审查；而是将这些表达作为一个整体，在其全部和完整的伯多禄传承中有效，正如我们常说的，并在适当之处将更详细地阐述。\par

% structure: p0012 source=h3 target=subsection reason=relative-heading-rank; base=h2

\subsection{教宗阿加托的书信}

\EditorialNote{见于Migne, Patrologia Latina, Tom. LXXXVII, col. 1161; 以及L. and C., Tom. VI, col. 630。}\par

\Signature{\Person{Agatho}——主教、天主众仆之仆——致最虔诚、安详的胜利者与征服者、我们至爱的、热爱天主和我们主耶稣基督的子孙、皇帝\Person{Constantine the Great}、\Person{Heraclius}与\Person{Tiberius}，奥古斯都们。}\par

正当我思虑人生各种忧虑，在唯一真天主前痛哭哀叹，祈祷祂将祂神圣仁慈的安慰赐予我摇摆的灵魂，并藉祂的右手将我高举出悲伤和焦虑的深渊时，我最显耀的主上们和儿子们，我极为感激地承认，你们举行会议的目的给了我深厚而奇妙的安慰。因为这是极其虔诚的，且源自你们最温和的平静，受神圣的良善所教导，为造福上主托付给你们守护的基督信仰国度，即你们的帝王权力和宽仁会关心殷勤探询有关天主的事（君王藉祂而统治，祂本身就是万王之王、万主之主），并寻求祂无瑕信仰的真理，正如宗徒们和宗徒时代的教父们所传下来的，并热切地吩咐在所有教会中持守纯洁的传统。为使无人不知你们这虔诚的意图，或怀疑我们是受强迫而非自由同意执行那些御旨——这些御旨关乎宣讲我们的福音信仰，是写给我们的前任、具宗徒记忆的教宗\Person{多努}的——它们已通过我们的职分被发送给并完全获得所有民族和人民的认可。因为这些御旨，圣神藉祂的恩宠，从纯洁心灵的宝库中，默启于皇帝执笔之舌，如一位顾问而非压迫者的话语，为自己辩护，而非轻视他人；不使人痛苦，而是劝勉；并以虔敬的方式邀请人归向那些属于天主的事。因为祂，所有人的创造者和救赎者，倘若以祂天主性的威严降临世界，本可惊吓凡人，却宁愿藉祂不可估量的仁慈和谦卑，下降至祂所创造的我们的境地，并以此救赎我们，祂也期望我们心甘情愿地宣认真信仰。\par

这正是真福伯多禄，宗徒之长，所教导的：\BibleQuote{你们要牧放在你们中间的基督的羊群，不是出于强迫，而是出于自愿，按照天主来劝勉他们。}因此，受这些谕旨的鼓励，万有之最温柔的主宰啊，并从痛苦的深渊中得到解脱，被提升到安慰的希望之中，我因一种更好的信心而稍有恢复，开始迅速遵行您最温柔坚毅的敕令先前所命令的事项，并努力依此寻找人员——如我们匮乏的时代及这顺服省份的素质所允许的——并与我的同仆主教们商议，既关系到这宗座即将召开的主教会议，也关系到我们自己的神职人员、基督宗教帝国的热爱者，之后又关系天主的虔敬仆人们，为劝勉他们迅速追随您最虔诚的安宁的足迹。并且，若不是我们谦卑的会议所在诸省的广大范围造成了如此巨大的时间损失，我们的卑微之身在不久之前本可以用勤勉的服从完成那即便现在也几乎未能完成的事。因为，当来自各省的会议正在我们周围聚集时，我们能够从直接隶属于您最宁静权力的这个罗马城本身，或从附近地区挑选一些人，另一些人我们又不得不从遥远省份等待，那些省份的基督信仰之言是由我卑微之人的前任们所派的人宣讲的；因此相当长的时间已经过去：我略过不提我身体的痛苦，由于这些痛苦，对于一个不断受苦的人来说，生命既不可能也不愉快。因此，最基督宗教的主宰和儿子们，依照您天主所保护的宽仁的最虔敬命令，我们以祈祷之心的虔诚（出于我们对您所欠的服从，并非因为我们倚靠我们所派去之人的[过度]知识），已注意派遣我们目前在场的同仆：\Person{Abundantius}、\Person{John} 和 \Person{John}，我们最可敬的兄弟主教们，\Person{Theodore} 和 \Person{George}，我们最亲爱的儿子和司铎，还有我们最亲爱的儿子 \Person{John}（执事），以及 \Person{Constantine}（这神圣精神母亲、宗座的副执事），以及 \Place{Ravenna} 圣教会的司铎代表 \Person{Theodore} 和天主的虔敬仆人们（隐修士）。因为，在被置于外邦人中、凭体力劳动相当分心地赚取日用粮的人们中间，怎么可能找到圣经的圆满知识呢？除非我们将我们圣洁与宗徒的诸位前任以及可敬的五大公会议所公义界定的，存于心灵的纯朴中，毫无歪曲地持守那来自教父们的信仰，始终渴望并努力拥有那唯一且至高的善，即：公义界定之事不可减少，不可更改，也不可添加，而应在文字和意义上原封不动地守护。对于这些同一专员，我们也给予了一些圣教父们的见证，这些圣教父是这基督的宗座教会所接纳的，连同他们的书籍；这样，他们既从您最慈祥的基督宗教的权柄获得了提议的特权，便可能借这些教父见证，努力提供满意答复（当您帝王的温柔如此命令时），论及这基督的宗座教会——他们精神上的母亲、您天主所生帝国的母亲——相信并宣讲什么，不是用世俗辞令的话——这是普通人所不能掌握的——而是用宗徒信仰的纯正，这种信仰我们从摇篮起就受教导，我们祈求直至末了都能服事并顺从上天之主、您基督宗教帝国的扩展者。因此，我们已授予他们能力或权柄，与您最宁静的威权一起，在您宽仁命令时以纯朴提供满意答复，但有限制：他们不得擅自添加、删除或改变任何事物；而应全然真诚地陈述这宗座的传统，正如我们前任的宗座教宗们所教导的。为这些代表，我们以心灵谦卑跪伏，恳求您始终充满屈尊俯就的宽仁，依照帝敕的最慈祥、最庄严的许诺，愿您如基督的安宁认为他们堪当接纳，并俯听他们最谦卑的建议。这样，愿您最温柔的热心找到全能天主的耳朵垂听您的祈祷，并愿您命令他们身怀我们宗座信仰的正直及身体的完整，安然返回自己的地方。同样，愿至高威权通过您天主所坚固的宽仁的最英勇、最不可征服的功绩，将整个基督宗教共同体恢复于您政府的仁慈统治之下，并愿他使敌国臣服于您大能的权杖，使从今以后每个灵魂和所有邦国都得到满足，因为关于那些来参加大公会议之人的豁免和安全，您曾通过您最庄严的书信郑重应许，您已在各方面实现。并非他们的智慧给了我们信心敢于派遣他们到您虔诚的御前；而是我们卑微之身顺从地履行了您帝王慈心以恩令所劝勉的。我们将简要向您这受天主教导的热心表明，我们宗座信仰的力量包含什么，这信仰我们是通过宗徒传统和宗座教宗们的传统，以及五大公圣会议的传统而领受的；藉着这些，基督的公教会的基础得以坚固和确立。这就是我们福音与宗徒信仰的状态[和正规传统]，即：正如我们宣认圣而不可分离的天主圣三，即圣父、圣子及圣神，是同一神性、同一本性、实体或本质，同样我们也将在它内宣认一个自然的意志、能力、行动、统治、威严、权能和光荣。并且，凡是就本质而言以单数形式论及同一圣三的，我们都理解为指向三位同质之位的同一本性，我们是由正典逻辑如此受教的。但是，当我们宣信有关那同一圣三的一位，即天主之子或天主圣言，以及祂按肉身的可敬救恩计划之奥迹时，我们断言，按照福音传统，所有事物在我们同一的主和救主耶稣基督内都是双重的，也就是说，我们宣认祂有两个本性，即神性和人性，从它们并在它们内，即使在奇妙而不可分离的结合之后，祂仍然存在。并且我们宣认，祂的每一本性有其自然的固有性，神性拥有神性的一切，毫无罪愆。并且我们承认，那同一降生成人——即成为人（\OriginalTerm{成为人的}{humanati}）——的天主圣言的每一（两个本性中的）本性，在祂内是不混淆、不可分离、不可改变的，惟独理智辨别出合一，以避免混淆的错误。因为我们同样憎恶分裂和混合的亵渎。因为当我们宣认我们唯一的主耶稣基督有两个本性、两个自然意志和两个自然行动时，我们并不断言它们是彼此相反或对立的（如同那些远离真理之路并指控宗徒传统的人所做的。愿这亵渎远离信友之心！），也不像在两个人位或位格中被分离（\OriginalTerm{本身分离}{per se separated}），而是说，正如我们同一的主耶稣基督有两个本性，同样祂也有两个自然意志和行动，即神性和人性：神性的意志和行动，祂从永远就与同质的圣父共有；人性的，祂从我们领受，在时间中与我们的本性结合。这就是宗徒和福音的传统，是您最幸福帝国的精神母亲——基督的宗座教会——所持守的。\par

这是虔诚的纯粹表达。这是基督宗教真实无瑕的信仰宣认，并非出于人的计谋，而是由圣神藉着宗徒之长所教导的。这是圣宗徒们坚固而无可指摘的教义，其中真诚虔诚的完整，只要被自由传扬，就捍卫你们陛下在基督化共和国的帝国，并使之欢欣（[将保卫它，使之稳固；并欢欣]），并且（我们坚信）将显示它充满幸福。请相信你们最卑微的（仆人），我最基督化的主上及儿子们，我正流着泪倾泻这些祈祷，为它的稳固与欢欣（\Emph{希腊文}作\InnerQuote{高举}）。这些事我（虽不配且微小）因真诚的爱而胆敢劝谏，因为你们天赐的胜利是我们的救恩，你们圣平的幸福是我们的喜乐，你们仁慈的无害是我们渺小的安全。因此，我以破碎的心和泪河、以俯伏的心神恳求你们，愿屈尊伸出最仁慈的右手，扶持那与你们虔诚劳苦的合作者、圣宗徒\Person{伯多禄}所交付的宗徒教义，使之不被藏在斗下，而是被吹号角更响亮地传遍全地：因为其真实的宣认，伯多禄因此被万有之主称为有福，是由天父所启示，因为他从万有救赎者亲口以三次托付接受了牧养教会灵性群羊的职责；在其保护盾下，他这宗徒教会从未偏离真理之路进入任何错误方向，其权威，作为诸宗徒之长的权威，整个天主教会及大公会议都忠信地拥抱，并在一切事上跟随；所有可敬的教父都拥抱其宗徒教义，藉此他们作为基督教会最被认可的明灯发光；圣洁的正统博士们都尊敬并跟随它，而异端者则以虚假控告和贬损仇恨追逼它。这是基督宗徒的活传统，他的教会到处持守，应当特别被爱戴和培育，并勇敢地传扬，它藉着真实的宣认使我们与天主和好，也使人在基督主前可称赞，它保护你们圣仁的基督化帝国，从天上之主赐给你们最虔诚圣勇以远播的胜利，它在战斗中陪伴你们，击败仇敌；它四面保护你们从天而生的帝国如同不可攻破的城墙，使敌对民族恐惧，以神圣义怒击打他们，在战争中从天赐予凯旋的棕榈枝用以打倒和征服敌人，并永远护卫你们最忠信的统治在平安中稳固快乐。因为这是真信仰的准则，你们最平静帝国的属灵母亲——基督的宗徒教会——无论顺逆都一直大力持守和捍卫；她，全能天主的恩宠将证明，从未偏离宗徒传统的道路，也未曾因屈服于异端的新奇而败坏，而是从起初从她的奠基者——基督的宗徒之长——接受了基督信仰，并保持无玷直到终结，按照主和救主自己的神圣许诺，他在神圣福音中对他的门徒之长说：\BibleQuote{伯多禄，伯多禄，看，撒殚想要得到你们，好筛你们像筛麦子一样；但我已为你祈求，叫你的信德不致丧失。待你回头以后，要坚定你的弟兄。}因此，请你们平静圣仁考虑，既然它是万有的主和救主，是他的信仰，他许诺伯多禄的信德不致丧失并劝他坚固他的弟兄，那么众所周知，宗徒教宗的先辈们，鄙人的前任，总是自信地做了这事：他们中，我们的卑微也藉着神圣指派接受了这职务，愿意成为跟随者，尽管不如他们且是众中最小的。因为我有祸了，若我忽略传扬我主的真理，而他们曾真诚传扬。我有祸了，若我以沉默掩盖真理，而这真理我受命交给兑换银钱的人，即教导基督徒子民并灌输给他们。在将来由基督亲自审查时，我若（天主不准！）在此为传扬他话语的真理而脸红，我将说什么？我当为自己，为托付给我的灵魂，能给出什么满意的交代，当他严格索要我已接受的职务之账时？那么，我最仁慈和最虔诚的主上及儿子们（我颤抖且心神俯伏说），谁不会被那奇妙的应许所激动，即对信者所说的：\BibleQuote{凡在人前承认我的，我在我天上的父前也必承认他}？又有哪个不信者不会被那最严厉的威胁所惊吓，他在其中充满义怒地抗议并声明：\BibleQuote{凡在人前否认我的，我在我天上的父前也必否认他}？因此，圣\Person{保禄}，外邦人的宗徒，也警告并说：\BibleQuote{但无论是我们，或是从天上来的天使，若给你们宣讲的福音与我们先前所宣讲的不同，当受诅咒。}既然如此，如此极端的惩罚悬在真理的败坏者或沉默的压制者头上，谁会不逃避对主信仰真理的任何削减呢？因此，鄙人的宗徒记忆的先辈们，在主的教义上有造诣，自从\Place{君士坦丁堡}教会的主教们试图将异端的新奇引入基督的无玷教会以来，从未停止以许多祈祷劝诫和警告他们，至少以沉默停止邪恶教义的异端错误，免得从此在教会合一中造成分裂的开始，主张在唯一耶稣基督我们的主内有两个本性的一个意志和一个运作：这是亚略派、阿波利纳里派、欧提奇派、蒂莫西派、无首派、特奥多西派和加伊亚尼派所教导的，以及一切异端的疯狂，无论是混淆基督降生成人奥迹的人，还是分裂它的人。那些混淆圣降生成人奥迹的人，因他们说基督的天主性和人性是一个本性，主张他有一个意志，如同一个（本性）和一个位格的运作。但那些分裂不可分结合的人，另一方面，将救主拥有的两个本性合起来，却不是在一个被承认为位格性的结合中；而是亵渎地通过意志的情感如同两个自立体即两个个体结合。此外，基督的宗徒教会，你们天立帝国的属灵母亲，宣认唯一耶稣基督我们的主存在于两个本性中并出自两个本性，并主张他的两个本性，即天主性和人性，在其不可分结合之后仍无混淆地存在于他内，并承认基督的每个本性的本性属性都是圆满的，并宣认属于本性属性的一切都是双重的，因为同一个我们的主耶稣基督自己既是完全的天主又是完全的人，出于两个并存在于两个本性中：在他奇妙的降生成人之后，他的天主性不可脱离他的人性而思考，他的人性亦不可脱离他的天主性。因此，按照圣而公的基督宗徒教会的准则，她也宣认并传扬他内有二个本性的意志和二个本性的运作。因为若有人意指位格意志，当在圣三中说有三个位格时，就必须主张三个位格的意志和三个位格的运作（这是荒谬且真正亵渎的）。因为正如基督信仰的真理所持，意志是本性的，当说到圣而不可分的天主圣三的一本性时，必须一致理解为有一个本性的意志和一个本性的运作。但当我们真实地宣认在我们的主耶稣基督——天主与人之间的中保——的独一位格内，有两个本性（即天主性和人性），甚至在他奇妙结合之后，正如我们正统地宣认同一独一位格的两个本性，同样我们也宣认他的两个本性的意志和两个本性的运作。为使这一真实宣认的理解从旧约和新约的天主所启示的教义中变得清晰，对你们虔诚的心（因为你们圣仁比我们卑微更能无可比拟地洞悉圣经的意义），我们的主耶稣基督自己，既是真实而完全的天主，又是真实而完全的人，在他的圣福音中，有时显示属人的事，有时显示属神的事，有时两者一起，作出关于他自己的显示，为教导他的信者相信并传扬他是真天主和真人。因此，作为人，他向父祈祷拿走苦杯，因为在他内我们的人性是完全的，除了罪以外：\BibleQuote{父啊！若是可能，就让这杯离开我吧！但不要照我所愿意的，而要照你所愿意的。}在另一处：\BibleQuote{不要照我的意愿，而要照你的意愿成就。}如果我们想知道这些证言由圣且被认可教父们解释的意义，并真正理解\Quote{我的意愿}和\Quote{你的意愿}意指什么，圣\Person{盎博罗削}在致有福记忆的\Person{格拉提安}皇帝的\WorkTitle{第二卷书}中用这些话教导我们这段的意义，说：\Quote{那么，他接受我的意愿，他接受我的忧愁，我自信地称它为忧愁，因为我是在谈论十字架，我的意愿是他称为他的，因为他作为人承担我的忧愁，他作为人说话，因此他说：\InnerQuote{不要照我所愿意的，而要照你所愿意的。}}我的悲伤是他按我的情感所接受的。看，最虔诚的君王，这位圣教父多么清楚地将主在祈祷中所用的话\Quote{不要照我的意愿}归于他的人性；藉此，按照圣宗徒\Person{保禄}的教导，他也被认为是\BibleQuote{服从至死，且死在十字架上}。因此我们也受教导，他服从他的父母，这应虔诚地理解为指他自愿的服从，不是按他的天主性（藉此他治理万有），而是按他的人性，藉此他自愿服从于父母。圣史\Person{路加}也见证同样的事，讲述同一的我们的主耶稣基督按人性向父祈祷并说：\BibleQuote{父啊！若是可能，就让这杯离开我吧！但不要照我的意愿，而要照你的意愿成就。}——\Person{阿塔纳修}，基督的宣信者、\Place{亚历山大里亚}教会的总主教，在其\WorkTitle{驳异端者阿波利纳里书}中，也理解为两个意志，如此解释：\Quote{当他说：\InnerQuote{父啊！若是可能，就让这杯离开我吧！但不要照我的意愿，而要照你的意愿成就。}又说：\InnerQuote{心神固然切愿，但肉体却软弱。}他显示有两个意志，一个是人性的即肉体的意志，另一个是属神的。因为他人性的意愿出于肉体的软弱逃避苦难，但属神的意愿为它准备着。}还能找到更真实的解释吗？因为在他内，即使不可分结合之后，按照大公会议的定义，有两个本性，怎能不承认他有两个意志，即人性的和属神的呢？因为\Person{若望}，那倚在主怀中的爱徒，也以这些话显示同样的自制：\BibleQuote{我自天降下，不是为执行我的旨意，而是为执行派遣我来者的旨意。}又说：\BibleQuote{派遣我来者的旨意就是：凡他交给我的，叫我一个也不失掉，而且在末日还要使他复活。}他又介绍主与犹太人辩论，在别的话中说：\BibleQuote{我不寻求我的旨意，而寻求那派遣我来者的旨意。}论及这些神圣话语的意义，圣\Person{奥斯定}在驳\Person{亚略派}的\Person{马克西米诺}的\WorkTitle{驳亚略派马克西米诺书}中如此写道。他说：\Quote{当子对父说\InnerQuote{不要照我所愿意的，而要照你所愿意的}时，你将自己的话置于服从之下并说：它显示他的意志服从了父，这对你有什么益处呢？好像我们会否认人的意志应服从天主的旨意似的？因为主是在他的人的性内说这话，任何人若仔细研读这段福音，会很快看出。因为在那里他说：\InnerQuote{我的心灵忧闷得要死。}这话能用于圣言的性体吗？但是，人啊，你试图使圣神的性体叹息，为何你说天主独生子的性体不能悲伤？但为防止有人这样辩论，他没有说\InnerQuote{我悲伤}（即使他这样说，也只能被理解为指他的人性），而是说\InnerQuote{我的心灵忧闷}，这是他作为人所拥有的心灵；然而在他所说的\InnerQuote{不要照我所愿意的}中也显示他愿意了父所不愿的某事，这只能在他的人性内发生，因为他没有将我们的软弱引入他的神性，而是要改变人性的情感。因为若他没有成为人，圣言绝不能对父说\InnerQuote{不要照我所愿意的}。因为那不变的性体绝不可能愿意不同于父所愿的。你们若作此区分，啊，亚略派，你们就不会是异端了。}\par

在这一论辩中，这位可敬的教父表明，当主说\Quote{自己的}时，他指的是他人性的意愿；而当他说不实行\Quote{自己的意愿}时，他教导我们不要主要寻求自己的意愿，而应通过服从将我们的意愿顺服于天主的旨意。由此一切显而易见，他拥有一个人性的意愿，借此服从他的圣父，并且他自身拥有这同一人性意愿，完全无罪，作为真天主和真人。圣盎博罗削在解释圣路加福音时也如此论述此事。\par

\Emph{此后是一系列教父引文，我认为不值得全部列出。在圣盎博罗削之后，他引用圣良，然后是圣额我略·纳齐盎，再然后是圣奥斯定。}（\WorkTitle{L. \& C.}，第647栏）\par

从这些证据可以清楚看出，属灵博士在此列举的每一种本性都有其自身的自然属性，且每一种本性都应被赋予一个意愿。因为天使的本性不可能拥有神圣的或人性的意愿，人性的本性也不可能拥有神圣的或天使的意愿。因为没有任何本性能够拥有属于另一种本性的任何事物或运动，而只能拥有受造时自然赋予的。既然这是事情的真相，那么最确定无疑的是，我们必须承认在我们的主耶稣基督内有两个本性和实体，即神性和人性，在其一个位格或位格中结合，并且我们进一步承认他内有二个自然的意愿，即神性的和人性的，因为他的神性就其本性而言不能被认为拥有人的意愿，他的人性也不应被认为自然地拥有神性的意愿。再者，基督的这两种实体中任何一种都不能被承认是没有自然意愿的；但他的人性意愿被其神性的全能所提升，他的神性意愿透过人性向人显示。因此，必须将属于神性的事归于他作为天主，将属于人性的事归于他作为人；并且透过同一位我们的主耶稣基督的位格性结合，每一样都必须被真实地承认，正如加采东大公会议最真实的法令所陈述的——\EditorialNote{此处为引文。}同一件事，在可敬的皇帝犹斯提念时代于君士坦丁堡召开的圣主教会议，在其定义的第七章中也教导。\EditorialNote{此处为引文。}此外，我们必须忠心地持守那些可敬的主教会议所教导的，以便我们永不因结合而消除本性的区别，而应承认基督为一位，真天主并真真人，每一本性的特性保持不变。因此，如果我们主耶稣基督的本性区别丝毫没有消除，那么我们必须在其一切特性中保存这同一区别。因为凡教导区别丝毫没有消除的人，宣告必须在一切事物中保存区别。但是，当异端者及异端的追随者说只有一个意愿和一种行动时，这区别又如何被认识？或这被本圣主教会议所定义的区别又保存在哪里？而如果断言他内只有一个意愿（这是荒谬的），那么作此断言的人必然要说那意愿要么是人的，要么是神的，要么是由两者混合、混淆而成的复合体，要么（按所有异端者的教导）基督只有一个意愿和一种行动，来自他（按他们所认为的）一个复合的本性。如此，本性的区别无疑就被摧毁了，而圣主教会议宣告即使在奇妙的结合之后也必须在一切方面保持它。因为，虽然他们教导基督是一位，其位格和实体是一位，但由于那由位格性结合而成的本性的结合，他们同样下令，在奇妙的结合之后，我们应当清楚承认并教导那在他内联合的本性的区别。因此，如果在同一位我们的主耶稣基督内，因（本性的）区别而保存了本性的特性，那么我们理应怀着圆满的信德也承认其自然意愿和行动的区别，以表明我们在一切方面都遵循了他们的教导，并且不接受任何异端的新奇事物进入基督的教会。\par

虽然其他诸位圣教父的著作甚多，但为免繁琐，我们仍从希腊文书籍中附上少许证据，附于我们这篇谦卑的陈述之后。\par

\Emph{此后是来自希腊教父的一系列引文，即：圣额我略神学家、圣额我略·尼撒、圣若望（君士坦丁堡主教）、圣济利禄（亚历山大主教）。}（\WorkTitle{L. \& C.}，第654栏）\par

从这些真实的证据也证明了，这些可敬的教父在同一位主耶稣基督内断言了两个自然的意愿，即神性的和人性的。因为当圣额我略·纳齐盎说：\Quote{那被理解为救世者之人的意愿}，他表明救世者的人性意愿因与圣言的结合而被神化，因此它不与天主相反。同样，他证明他拥有一个虽是人性却被神化的意愿，并且这同一个意愿，他（如前所述）也拥有，与他的神性意愿一同，而神性意愿与父的意愿是同一的。因此，如果他拥有一个神性的和一个被神化的意愿，他也就有两个意愿。因为本性为神性的不需要被神化；被神化的不是本性上真正的神性。而当伟大的主教圣额我略·尼撒说，这奥迹的真正宣认是应当理解基督内有一个人的意愿和另一个神性的意愿，他说一个和另一个意愿，除了显然有两个意愿之外，还叫我们理解什么呢？\par

\Emph{他接着评论引自圣若望的段落，然后评论引自亚历山大圣济利禄的段落。此后依次是引自圣希拉利、圣亚大纳削、圣狄奥尼修（亚略巴古人）、圣盎博罗削、圣良、圣额我略·尼撒、亚历山大圣济利禄的引文，并按顺序加以评论。他然后继续写道：}（\WorkTitle{L. \& C.}，第662栏）\par

在其他可敬的教父们那里不乏最有力的段落，他们清楚地讲论基督内的两个本然行动，更不用说圣济利禄·耶路撒冷、圣若望·君士坦丁堡，或后来为保卫可敬的加采东大公会议及圣良的《大卷》而对抗异端者（从他们的错误中产生了这个新教义的断言）的人：即斯基托波利斯的若望主教、亚历山大的欧洛吉主教、尤弗拉米乌斯和年长的阿纳斯塔修斯（提奥波利斯教会最配得上的统治者），以及超越所有的是那位真正宗徒信仰的效仿者、可敬记忆的查士丁尼皇帝，他的正直信仰高扬了基督国度，正如他真诚的宣认取悦了天主一样。他的可敬记忆被万民尊崇，他的正直信仰通过他极其庄严的谕令在整个世界被传扬赞美：其中之一，即致亚历山大的佐伊禄宗主教关于反对无首异端以满足他们对宗徒信仰正直性的谕令，我们随同这份我们卑微的文书，通过同样的信使，呈献给阁下最宁静的基督信仰。但为了不让这一陈述因其长度而显得繁琐，我们在这份谦卑的陈述中只插入了少数圣教父的见证，尤其是在写给那些被公认为整个世界如稳固基础般依赖其照顾和安排的人时；因为这是无可比拟和伟大的，即为了真宗教的热爱和热诚而暂时放下整个基督国度的关怀时，阁下至尊至虔诚的仁慈愿意更清楚地理解宣讲的教义。因为从不同的受人称赞的教父那里，正统信仰的真理已经变得清晰，尽管论述简短。因为受人称赞的教父们认为对明显清楚的事详加论述是多余的；因为即使愚钝的人，谁看不出明显的事呢？因为一个本性没有本然行动是不可能的，也是违反自然秩序的：甚至连异端者也不敢这样说，尽管他们都在为反对信仰的正统而寻求人的狡猾和诡诈问题，以及符合他们堕落的论点。\par

那么，如今如何能断言那圣善的正统教父从未说过、甚至亵渎的异端者也未曾胆敢发明的事，即：在基督的两个本性（神圣的和人性的，每一个的固有特性都被认为在基督中得以保持）中，任何心智健全的人竟宣称只有一个行动？因为如果只有一个，那就请他们说那是暂时的还是永恒的，神圣的还是人性的，受造的还是非受造的：与圣父相同的还是与圣父不同的？因此，如果只有一个，那同一个必须同时属于神性和人性（这是荒谬的），那么当既是天主又是人的天主子在地上行人性的事时，圣父也照样按其本性（\OriginalTerm{本然上}{naturaliter, \GreekText{φυσικῶς}}）与他同工；因为圣父所做的事，圣子也照样做。但如果（正如事实那样）基督所行的人性行为只应归于他作为圣子的位格（这与圣父的位格不同），那么在一个本性中基督行了一类事，在另一个本性中行了另一类事，以至于按照他的神性，圣子做与圣父同样的事；同样，按照他的人性，那些属于人性的合适行为，他也作为人行了，因为他确实是真天主真人。因此我们正确相信，同一个位格，既然他是独一的，就有两个本然行动，即神圣的和人性的，一个是非受造的，另一个是受造的，作为真而完美的天主和真而完美的人，那同一个、唯一的、天主与人之间的中保，主耶稣基督。因此，从行动的性质，我们认识到在基督内通过位格结合所联合的本性之间存在一种\OriginalTerm{无冒犯}{\GreekText{ἀπρόσκοπος}}的差异。我们现在引用一些来自令天主憎恶的异端者可憎著作的段落（我们同样憎恶他们的言论和话语），以证明那些我们新教义的发明者所追随的、教导在基督内只有一个意志和一个行动的东西。\par

\EditorialNote{随后引用阿波利纳里、塞维鲁、狄奥多西（亚历山大）的段落。（拉比与琴茨版，第667栏。）}\par

看，最虔诚的主们和儿子们，通过圣教父们的见证，如同通过属灵的光芒，公教与宗徒教会的教义已被阐明，异端盲目的黑暗（它正在向人提供谬误以供模仿）已被揭示。现在，这新教义必须追随某人，且是靠谁的权威支持，我们将指出。\par

\EditorialNote{此处引用亚历山大的居鲁士、法龙的狄奥多、君士坦丁堡的塞尔吉、彼鲁斯、其继任者保禄、其继任者伯多禄的段落。（拉比与琴茨版，第670栏。）}\par

那么，愿基于天主的仁慈，以内在的辨别之眼（您因天主的恩宠而被认为配得接受用于引导基督徒人民），请注意您认为基督徒人民应当追随哪一位这样的博士，应当拥抱他们中哪一位的教义以得救；因为他们彼此谴责，每个人谴责另一个人，按照他们著作中各种不稳定的定义，有时断言有一个意志和一个行动，有时断言既没有一个也没有两个行动，有时一个意志和行动，又两个意志和两个行动，同样一个意志和一个行动，又再次既不是一个也不是两个，而另一个人说一个和两个。\par

谁若不渴望得救，并寻求在主来临时献上正直的信德，岂能不憎恨、愤怒并回避如此盲目的谬误？因此，靠着您的全力（藉天主的助佑），天主圣教会——您最基督徒权能的母亲——应从这类教师的谬误中被解救和解放；而正统信仰的福音和宗徒的正直，已建立在这蒙福的\Person{圣伯多禄}（宗徒之长）的教会的坚固磐石上，并藉他的恩宠和护佑免于一切谬误；我所说的那信仰，全体统治者、司铎、圣职人员和人民，都应与我们一起一致宣认并宣讲，作为宗徒传统的真实宣言，为悦乐天主并拯救自己的灵魂。\par

我们辛苦将这些事插入我们卑微的论著中，因为我们一直痛苦并不断叹息，如此严重的谬误竟被教会的主教们接纳，他们热衷于建立自己的特殊观点而非信仰的真理，并认为我们真诚的弟兄劝告源于对他们的轻蔑。事实上，我卑微的宗座前任们曾劝诫、恳求、指责、哀求、斥责，并运用各种劝勉，使这新伤口得以医治；他们做出此举并非出于仇恨之心（天主为我作证），也非出于夸耀的得意、争辩的对立、或空虚地希望寻找他们教导的缺点，亦非类似傲慢之爱，而是出于对真理正直的热忱、对纯正福音宣认的规范、对灵魂的得救、对基督徒国家的稳定、以及对罗马帝国统治者的安全。他们也未因这长期久已根深蒂固的谬误而停止劝诫，而是一直劝勉并作证，且出于弟兄之爱，并非出于恶意或固执的仇恨（远离基督徒之心吧，不要因他人跌倒而喜乐，万有的主教导说：\Quote{我不愿罪人死亡，而愿他悔改得生}；而谁为那一个悔改的罪人欢喜，超过为那九十九个义人呢？他由天降下到地上，为拯救迷失的羊，屈尊了他威能的能力），而是伸开属灵的臂膀渴望他们，劝勉他们拥抱回归正统信仰的合一，并期待他们悔改至正统信仰的完美正直：使他们不致自绝于我们的共融，即与宗徒\Person{圣伯多禄}的共融——我们（虽不配）执行其职务，并宣讲他所传授的信仰——而应与我们一同向基督主——无瑕的祭献——祈祷，为你们最强大和平安的帝国之稳固。\par

我们相信，最虔诚的万有之主（拉丁文为单数），已不存在任何可能妨碍认出那些跟随新教义发明者之人的模糊之处。因为教父们的言论充满属灵理解的甘甜，如今已显明于众人眼前；而异端者的臭气，所有信徒都当躲避，也已众所周知。同样也非秘密，新教义的发明者已被证明是异端者的追随者，而非步武圣教父们足迹的人：因此，凡企图粉饰自己任何谬误的人，都被真理之光照耀定罪，正如外邦人的宗徒所说：\Quote{凡彰显的即是光}，因为真理常是恒定不变的，而谬误则常变化，在其游荡中采取相互矛盾之事。为此，新教义的发明者已被证明教导了相互矛盾的事，因为他们不愿追随福音和宗徒的信仰。因此，既然真理藉着你们怀有天主所赋虔诚的观察而照耀出来，已被揭露的谬误也获得了应有的轻蔑，剩下的就是使加冕的真理藉着你们蒙天主加冕的仁慈之虔诚眷顾而胜利照耀；并且使新奇的谬误及其发明者和跟随其学说的人，因他们的傲慢而受到应得的惩罚，并从正统主教们中被驱逐出去，因为他们异端邪恶的新奇事物，他们曾试图引入圣而公的宗徒的基督教会，并以异端邪恶的传染玷污基督教会不可分和无瑕的身体。因为有害者损害无辜者是不公义的，也不该让一些人的过犯临到无过之人身上；因为即使在这世上，对受定罪者施以慈悲，但那些如此被宽恕的人对于那次宽恕在天主的审判中得不到任何益处，而那些如此宽恕他们的人则因他们不合法的怜悯而招致不小的危险。\par

但是，我们相信，全能的天主已将纠正这些事，保留给你们仁爱之主的幸福日子，为使你们在地上充满我们的主耶稣基督亲自的地位和热忱——祂曾恩准加冕你们的统治——使你们能为祂的福音和宗徒的真理作出公正的审判。因为，尽管祂是人类的救赎者和救主，却受了伤害，并且一直忍受至今，感动了你们刚毅的帝国，使你们值得追随祂信德的事业（正如公平所要求，并如圣教父和五次大公会议所断定的），并且你们应当在祂的护佑下，为那些蔑视祂信德的人，对你们的救赎者和共同为王所受到的伤害施行报复，以此以君王的仁慈宽宏大量地实现了那个先知性的言语，就是达味王和先知向天主所说的：\BibleQuote{我为你殿宇的热忱把我耗尽。}因此，因这如此悦乐天主的熱忱而被称赞，他理当听到由众人的造物主所说的这句有福的话：\BibleQuote{我找到了达味，一个合乎我心的人，他要承行我的一切旨意。}并且在圣咏中对他也有应许：\BibleQuote{我找到了我的仆人达味，以我的圣油给他傅油：我的手要扶持他，我的臂膀要坚强他。}如此，你们基督的仁爱之最虔诚的陛下，可以为了报酬而用炽热的熱忱来推进基督的事业，并且愿那在圣福音中储存了应许的那一位，使你们最强盛帝国的一切行为都幸福而兴盛，祂说：\BibleQuote{你们先寻求天主的国，这一切自会加给你们。}因为所有知道这神圣纲领的人，已经向你们最强盛帝国的护卫者献上了无数的感恩和不尽的赞颂，充满了对你们仁爱之伟大的钦佩，因为你们如此仁慈地显明了你们庄严宽大的善意；实在，作为最虔诚和最公义的君主，你们已屈尊以敬畏天主来处理神圣的事，并已承诺给予从我们卑微处派往你们那里的那些人一切豁免权。\par

并且我们确信，你们虔诚的仁爱所应许的，你们有能力实行，为使那由超越你们基督徒能力的宗教爱德向天主所誓许和应许的，仍能藉祂的全能得以成就。\par

因此，愿所有基督民族，以及永恒的纪念和不断的祈祷，在主基督面前倾流，因为这是祂的事业，为你们的平安、你们的胜利和你们完全的凯旋；为使外邦人各民族，为至高的威严所震慑，最谦卑地将他们的颈项俯伏在你们最强大统治的权杖之下，为使你们最虔诚国度的权力得以持续，直到永恒国度无尽的喜乐接替这暂时的统治。再也没有什么比这更能将你们不可战胜的刚毅之仁爱，推荐给天主的威严，那就是使那些偏离真理准则的人被驳斥，并使我们的福音和宗徒信仰的完整，得以到处宣讲和传布。\par

此外，最虔诚、受天主训导的儿子与主上，倘若君士坦丁堡教会的总主教愿意与我们一同坚持并宣讲圣经、可敬的会议和神性教父们这无可指责的宗徒道理法则，依循他们对福音的理解，我们藉着圣神的助佑已将真理的形式借此阐明，那么，爱慕天主圣名的人必得享大平安，分裂的绊脚石必不复存在，并且\BibleBook{}{宗徒大事录}所记载的事必将实现：当人民藉着圣神的恩宠认识基督信仰时，我们众人将同心合意。但是，倘若他（愿天主禁止！）竟然偏爱接纳他人最近引入的新奇事物，并以与正统真理法则及我们宗徒信仰相悖的道理自陷网罗；而为了规避这些有损灵魂的谬论，他们不顾宗座先辈们的劝诫和警告，至今仍有所延迟，那么他本人应当知道，在基督的神圣审判中，在全能者面前，当基督来审判时，我们也必须为我们所负的宣讲真理的职务交账，或为我们容忍违背基督信仰的事物交账，他将为如此藐视付出怎样的答复。但愿我们（我谦卑地祈求）能不受混淆、自由地、以纯朴和纯洁、完整无缺地持守我们从起初所接受的正统信仰的宗徒与福音法则。愿陛下至为庄严的安宁，因您对公教与宗徒正统信仰的挚爱与尊敬，从我们的主耶稣基督——与您一同治理您基督帝国的统治者——获得您虔诚劳苦的圆满赏报，您渴望保持祂的真实宣认无玷，因为您那受天主加冕的仁慈丝毫没有忽略或遗漏任何有助于教会平安的事，只要始终维护真正信仰的完整：因为天主，万民的审判者，按祂认为最适宜的安排万事的结局，祂洞察内心的意图，并悦纳虔敬的热忱。因此，我劝勉您，最虔诚、宽仁的皇帝，连同我这卑微之人，每一位基督徒都屈膝谦卑地劝勉您，为了一切悦乐天主的善行和令人钦佩的帝国恩惠——天主的屈尊俯就藉着您蒙天主悦纳的照料赐予人类——您也要下令，为了圆满虔敬的重新整合，向与您共治的基督主奉上可悦纳的祭献，赐予完全豁免权，并给予每位愿意发言者自由表达的权利，为他自己所信仰和持守的道理辩护，如此最明显地让众人知晓，没有人因恐惧、暴力、威胁或厌恶而禁止或阻止任何愿意为公教与宗徒信仰真理发言的人，并且所有人心同意和地颂扬您\OriginalTerm{神圣的}{divinam}威严，一生为如此伟大、不可估量的恩惠，并不断向基督主献上祈祷，愿您最强大的帝国得以保全、不受损害、得到高举。\Signature{愿来自上天的恩宠保全您的帝国，最虔诚的主上，并将万国之颈置于其足下。}\par

% structure: p0038 source=h2 target=section reason=relative-heading-rank; base=h2

\section{\Person{阿加托}与一百二十五位主教之罗马会议的函件，作为出席第六届大公会议之代表团的指示。}

（见《\Person{拉比}与\Person{科萨特}编》\WorkTitle{公会议}，第VI卷，第677栏及以下；并见《\Person{米涅}编》\WorkTitle{拉丁教父集}，第LXXXVII卷，第1215栏及以下。[我遵循的最后一个文本是\Person{曼西}的版本。]）\par

\Signature{致最虔诚的主上、最宁静的胜利者和征服者，我们蒙天主和我们的主耶稣基督所爱的儿子，君士坦丁大帝，以及赫拉克利乌斯和提比略，奥古斯都，\Person{阿加托}，主教暨天主众仆之仆，连同隶属于宗座会议的所有会议。}\par

\EditorialNote{该信函以对皇帝的多项赞美开头，其风格和内容与前信引言相似。笔者认为不值得翻译，而直接从教义部分开始，其全文呈现给读者。（拉比与科萨特，第682栏。）}\par

我们信全能的天主圣父，天地万物，无论有形无形，都是他所创造的；又信他的独生子，在万世之前由他所生，真天主由真天主，光由光，受生而非受造，与圣父同体；万物由他而造；又信圣神，主及赋予生命者，由圣父所发，与圣父圣子同受钦崇，同享光荣；天主圣三在合一中，合一在天主圣三中；就本质而言是合一，但就位格或自立存在而言是三；因此我们宣认天主圣父、天主圣子、天主圣神；并非三神，而是一神，圣父、圣子、圣神；不是三个名称的自立存在，而是三个自立存在的同一实体；在这些位格中有一个本质、或实体、或本性，也就是说，有一个神性、一个永恒、一个权能、一个王国、一个荣耀、一个钦崇、一个本质性的意志和行动，属于同一圣而不可分的天主圣三，他创造了万物，安排万物，并仍保持万物。\par

再者，我们明认同一圣洁、同体之圣三中的一位，天主圣言，祂在万世之前由圣父所生，在末世为了我们并为了我们的救恩，自天降下，由圣神及我们的元后——圣洁、无玷、永远童贞、光荣的玛利亚，真实而确切的天主之母——降生成人，即按由她所生的肉身而论；祂真实地成为了人，同一者是真天主亦是真人。祂出自天主圣父是天主，但出自童贞母亲是人，由她的肉身降生成人，具有理性灵魂；按天主性言，与天主圣父同体；按人性言，与我们同体，在一切事上相似我们，只是没有罪。祂在般雀·比拉多下为我们被钉十字架，受难，被埋葬，复活了；升了天，坐在圣父的右边；祂将再来审判生者死者，祂的王国没有终结。\par

我们承认我们同一的主耶稣基督、天主的独生子，存在于并出于两个性体，无混淆、无变化、无分裂、无分离；结合并未消除性体间的差异，反而保留了每一性体的特性，并合于一位格、一个存在，并非分裂或划分为两个位格，亦非混淆为一个复合性体；我们明认同一独生子、天主圣言、我们的主耶稣基督，并非一个在另一个之内，也非一个附加于另一个，而是祂自己在两个性体之内——即在神性与人性之内，即使在位格结合之后亦然：因为圣言并未改变为肉身的性体，肉身亦未转化为圣言的性体，因每一性体都保持了其本性。我们仅凭默观分辨在祂内结合的各性体之间的区别，祂由这些性体无混淆、无分离、无变化地构成；因为二者合而为一，并借着一个二者同在，因为神性之崇高与肉身之卑微同时并存，每一性体在结合后保持其本有特性而无任何缺失；每一\Quote{形式}在与另一\Quote{形式}的共融中行其所特有之事。圣言行圣言之事，肉身行肉身之事；前者以奇迹发光，后者以凌辱弯腰。因此，正如我们明认祂实有两个性体或本体，即神性与人性，无混淆、无分离、无变化地结合，同样，虔诚的规则教导我们，祂有两个自然的意志和两个自然的运作，作为完美的天主和完美的人，同一的我们的主耶稣基督。宗徒与福音的传承以及圣教父（圣宗徒及大公教会与可敬的各届大公会议所接受者）的权威已明确教导了我们。\par

\EditorialNote{信函继续说明这是传统信仰，也是教宗玛尔定所主持的会议所阐明的；反对这信仰的人已偏离真理，有的这样，有的那样。接着为延迟派遣皇帝敕令所指定的人员而道歉，并如此继续：（拉贝与科萨，第686栏；米涅，第1224栏）。}\par

首先，我们中许多人分散在广大的地域，直至海滨，他们的旅程必然花费许多时间。此外，我们曾希望能将我们的同仆及兄弟主教、大不列颠岛的总主教及哲学家戴奥多若，连同迄今仍被留在那里的其他人，联合于我们的卑微之中；并在这会议中增添各地教区的不同主教，以便我们卑微的建议（即信函中所含的教义定义）能来自一个影响广泛的会议，以免只有一部分人知情而另一部分人不知；尤其在外邦人中，如伦巴第人、斯拉夫人，以及法兰克人、法国人、哥特人、不列颠人，我们知道有许多同仆不断好奇地探问此事，以了解在宗徒信仰的事业中正在做什么：他们只要与我们持有同一真信仰、思想一致，便能有益；但如果（愿天主禁止！）他们在任何信德条款上绊倒，就会成为麻烦和对立。然而我们，虽然最为卑微，仍竭尽全力使你们基督教帝国的国祚得以显示超越万民，因为在其中建立了真福伯多禄——宗徒之长——的圣座，其权威使所有基督教民族因着对真福宗徒伯多禄本人的尊敬，与我们一同敬拜和尊崇。（这是拉丁文，在我看来有讹误；希腊文如下：\Quote{其权威为真理之故，所有基督教民族与我们一同敬拜尊崇，依照对真福伯多禄宗徒本人的尊敬。}）\par

\EditorialNote{信函以祈求恒心、为国家和皇帝祝福、并期望真理普世传播与接受而结束。}\par

% structure: p0048 source=h2 target=section reason=relative-heading-rank; base=h2

\section{第八次会议——会议记录摘录}

\EditorialNote{皇帝说}\par

让格奥尔基，我们这由天主保佑之城的最神圣总主教，以及玛加略，安提约基雅的可敬总主教，以及隶属于他们的主教会议（即他们的附属主教）说，他们是否服从由最圣洁的阿加托，古罗马教宗及其主教会议所送建议的\OriginalTerm{效力}{\GreekText{εἰ} \GreekText{στοιχοῦσι} \GreekText{τῇ} \GreekText{δυνάμει}}。\par

\EditorialNote{格奥尔基的回答，他的所有主教（许多主教逐一发言）都同意，除了梅蒂莱内的戴奥多若（他在第十次会议结束时提交了同意书）。}\par

我曾勤勉地审查了呈给陛下至敬虔诚的建议书中的所有论据，这些建议书是由古罗马至圣教宗\Person{阿加托}及其会议所呈，我也仔细核查了存放在我尊敬的宗主教区内的圣教父及被认可的教父们的著作。我发现，那些建议书中所包含的所有圣教父及被认可的教父们的证言，与圣教父及被认可的教父们完全一致，并无任何差异。因此，我服从他们，并如此宣认和相信。\par

[\Emph{所有其他隶属君士坦丁堡宗座的主教的回答}。（第735栏。）]\par

我等，至敬之主，接受由古罗马至圣可敬教宗\Person{阿加托}及其属下的会议所呈给陛下至仁毅勇的建议书的教导，并遵循其中所含的意旨，如此我等思虑，如此我等宣认，且如此我等相信：在我等唯一主耶稣基督、我等真天主内，有两个本性，不混乱、不变换、不分离，以及两个自然意志和两个自然运作；凡曾教导或如今说在我等唯一主耶稣基督、我等真天主的两个本性内只有一个意志和一个运作的人，我等一概予以绝罚。\par

\Emph{皇帝对玛卡里的要求}。（第739栏。）\par

让安提约基雅可敬的总主教\Person{玛卡里}，他已听到了这神圣公会议所说的话〔要求他表明信仰〕，回答他认为合适的话。\par

\Emph{玛卡里的回答}。\par

我不说在我主耶稣基督的降生奥迹中有两个意志或两个运作，而是一个意志和一个\OriginalTerm{神人合一之运作}{theandric operation}。\par

% structure: p0059 source=h2 target=section reason=relative-heading-rank; base=h2

\section{第十三场会议——针对一志论者的判决}

神圣会议说：我等根据向陛下所作的承诺，重新审查了先前曾为这上天保佑的皇城的宗主教\Person{塞尔吉乌斯}致当时是法息斯主教的\Person{基禄}的信，以及致曾任古罗马教宗的\Person{霍诺里乌斯}的信，还有后者给同一\Person{塞尔吉乌斯}的回信。我等发现这些文件与宗徒信理完全无关，与神圣公会议的宣言以及所有被认可的教父们无关，它们追随了异端者的错误教导；因此我等完全拒绝它们，并视之为害及灵魂而予以憎恶。但那些我等憎恶其学说之人的名字也必须从神圣天主的教会中逐出，即：先前曾为这上天保佑之城的宗主教\Person{塞尔吉乌斯}，他是第一个写这邪恶学说的人；还有\Person{基禄}（亚历山大里亚的）、\Person{皮尔赫}、\Person{保禄}和\Person{伯多禄}，他们都在去世时任这上天保佑之城的主教，并与他们心志相同；还有曾任法君主教的\Person{德奥多禄}。所有这些人均已被古罗马至圣三度可敬教宗\Person{阿加托}在其致我等至敬虔诚的上主及大能皇帝的呈文中拒绝，因为他们违背我等正统信仰的心志；我等规定他们均应受绝罚。此外，我等规定曾任古罗马教宗的\Person{霍诺里乌斯}也应从神圣天主的教会中逐出并受绝罚，因为我们在其致\Person{塞尔吉乌斯}的信中发现，他在各方面都追随\Person{塞尔吉乌斯}的观点并确认了他的邪恶学说。我等也审查了圣善记忆的\Person{索弗洛尼乌斯}（先前曾为基督我等真天主圣城耶路撒冷的宗主教）的公函，发现它符合真信仰、宗徒教导以及被认可的圣教父们的教导。因此我等接纳它为正统并有益于神圣公教会和宗徒教会，并规定他的名应载入圣教会的双连画中。\par

% structure: p0061 source=h2 target=section reason=relative-heading-rank; base=h2

\section{第十六场会议——会议记录摘录}

\Emph{教父们的欢呼}。\par

皇帝万岁！\Person{君士坦丁}，我伟大的皇帝，万岁！正统的君王万岁！使我等享太平的皇帝万岁！\Person{君士坦丁}，新\Person{马尔西安}，万岁！\Person{君士坦丁}，新\Person{德奥多西}，万岁！\Person{君士坦丁}，新\Person{犹斯提尼安}，万岁！正统信仰的守护者万岁！求主护佑教会的基础！求主护佑信仰的守护者！\par

罗马教宗\Person{阿加托}万岁！君士坦丁堡宗主教\Person{乔治}万岁！安提约基雅宗主教\Person{德奥法努斯}万岁！正统的公会议万岁！正统的元老院万岁！\par

对\Place{法兰}的\Person{德奥多禄}，异端者，绝罚！\newline{}对\Person{塞尔吉乌斯}，异端者，绝罚！\newline{}对\Person{基禄}，异端者，绝罚！\newline{}对\Person{霍诺里乌斯}，异端者，绝罚！\newline{}对\Person{皮尔赫}，异端者，绝罚！\newline{}对\Person{保禄}，异端者，绝罚！\newline{}对\Person{伯多禄}，异端者，绝罚！\newline{}对\Person{玛卡里}，异端者，绝罚！\newline{}对\Person{斯德望}，异端者，绝罚！\newline{}对\Person{波利克罗尼乌斯}，异端者，绝罚！\newline{}对\Place{培尔加}的\Person{阿佩尔吉乌斯}，异端者，绝罚！\newline{}对所有异端者，绝罚！对所有附和异端者的人，绝罚！\newline{}愿基督徒的信德增长，愿正统的公会议享寿长久！\par

% structure: p0066 source=h2 target=section reason=relative-heading-rank; base=h2

\section{信仰的定义}

\Emph{（见于会议记录，第十八场会议，L. and C.，\WorkTitle{公会议}，第六卷，第1019栏。）}\par

神圣、伟大的公会议，因天主的恩宠和至为虔诚、忠信及大能的\Person{君士坦丁}元首的宗教法令，在这上天保护和皇家的君士坦丁堡新城、新罗马，于皇宫大厅称为\Place{特鲁卢斯}者内，集会并规定如下。\par

天主圣父的独生子和圣言，他在一切事上同我们一样，只是没有罪，他成了人，我们真正的天主基督，在福音之言中明确宣告说：\Quote{我是世界的光；跟随我的，决不在黑暗中行走，必有生命的光。} 又说：\Quote{我把我的平安留给你们，我将我的平安赐给你们。} 我们最温和的君主，正统信仰的捍卫者、邪恶教义的反对者，因着这由天主发出的平安之言而虔敬地受引导，召集了这次神圣的大公会议，联合了整个教会的判断。因此，我们这次神圣的大公会议，驱散了迄今一度流行的不虔敬的错误，并且紧密追随圣教父们所认可的正直道路，虔敬地完全赞同五次神圣大公会议：即聚集在尼西亚反对狂暴的亚略的三百一十八位圣教父的大公会议；其次是在君士坦丁堡由一百五十位受天主感召的教父反对圣神的反对者马其顿尼和不虔敬的阿波利纳里的大公会议；又是在厄弗所由二百位可敬人士首次聚集反对犹太派者奈斯多利的大公会议；以及在加采东由六百三十位受天主感召的教父反对天主所憎恨的欧提基和狄奥斯库若的大公会议；此外，还有最近一次，即在本城聚集的第五次神圣大公会议，反对摩普绥提亚的狄奥多若、奥力振、狄迪慕斯、埃瓦格略，以及狄奥多莱反对著名的济利禄《十二章》的著作，和据说由依巴斯写给波斯人玛里斯的书信。我们全面更新了宗教的古代法令，驱散了不虔敬者的邪恶教义。我们这次受天主启示的神圣大公会议，批准了由三百一十八位教父颁布、后又由一百五十位教父虔诚确认的信经，其他神圣会议也衷心接受并批准了该信经，以铲除一切毁灭灵魂的异端。\par

三百一十八位圣教父的尼西亚信经。\par

我等信，等等。\par

聚集在君士坦丁堡的一百五十位圣教父的信经。我等信，等等。\par

神圣的大公会议进一步宣认，这虔敬且正统的信经出自天主的恩宠，本足以使人完整认识并确认正统信德。然而，那从起初藉着蛇为助、以此向人类投下死亡之毒的邪恶创始者，并未罢休；如今他同样找到了适于执行其旨意的工具——即费兰主教德奥多禄、这皇城之总主教塞尔吉、彼鲁斯、保禄、伯多禄，以及古罗马教宗和诺理、亚历山大里亚主教居鲁士、不久前就任安提约基雅主教的玛加略及其弟子斯德望——积极利用他们，为整个教会设置了绊脚石，即在我们真实的天主、圣三之一的基督的两个性体内只有一个意志、一个行动的谬论；他们用新造术语向正统信众散布一种异端，与不虔敬的阿波利纳里、塞维鲁和特米斯提乌斯的疯狂邪恶教义相似，并试图狡诈地摧毁我们主耶稣基督、我们的天主降生成人的圆满性，亵渎地宣称他那具有理性灵魂的肉身是无意志、无行动的。因此，我们的天主基督兴起了我们虔信的君主，一位新的达味，按经上所记载，他找到了一位合他心意的人，即\BibleQuote{不容自己的眼睛睡眠，也不容自己的眼睑打盹}，直到他藉着这由天主聚集的神圣大公会议找到了完美的正统宣言；因为按照天主说的话：\BibleQuote{哪里有两个或三个人因我的名聚在一起，我就在他们中间}。这神圣的大公会议欣然接受并举手致敬至圣可敬的古罗马教宗阿加托寄给我们至虔至信的君士坦丁皇帝的谏言，该谏言点名拒绝了那些在基督我们的天主降生成人的救世工程中教导或宣讲一个意志、一个行动的人，也采纳了同一至圣教宗领导下、由125位可敬的主教组成的会议寄给他的联名谏言，视其与神圣的加采东大公会议及同一古罗马教宗至圣可敬的良一世致圣弗拉维安的\WorkTitle{书函}一致——该\WorkTitle{书函}被本会议称为正统信德的支柱——同时也赞同真福济利禄反对不虔敬的聂斯多略、写给东方各主教的联名信函。追随五个神圣的大公会议以及神圣可敬的教父们，我们异口同声地界定：必须承认我们的主耶稣基督是真实的天主和真实的人，是圣而同一性体的、赋予生命的圣三之一，在天主性内圆满，在人性内也圆满，真实的天主和真实的人，由理性灵魂和人的身体所构成；按天主性而论与圣父同一性体，按人性而论与我们同一性体；除了罪以外，在一切事上与我们相似；按天主性而言，在万世之前由圣父所生，但在这末后的时日，为了我们人类并为了我们的救恩，由圣神和童贞玛利亚——严格而正确地称为按血肉的天主之母——取得人性；同一基督、主、独生子，具有两个性体，不混乱、不变、不可分、不离散，各性体的特性在结合中并未丧失，反而每一性体的属性得以保存，结合于一个位格、一个自立体，不被分离或分割为两个位格，而是同一的、独生的天主圣言、主耶稣基督，正如古先知们所教导的，以及我们的主耶稣基督亲自教导的，圣教父们的信经所传授的；界定这一切后，我们同样宣告：在他内有两个本性的意志和两个本性的行动，不可分、不变、不离散、不混乱，按圣教父们的教导。这两个本性的意志并非彼此对立（但愿没有！）——像不虔敬的异端者所断言的那样——而是他的人性意志跟随，并非抵抗或勉强，而是服从于他的天主性和全能意志。因为按照最智慧的亚大纳削所说，肉身本当被推动而服从天主的意志。正如他的肉身被称为且确实是天主圣言的肉身，同样地，他肉身的本性意志也被称为且确实是天主圣言的本有意志，正如他亲自说的：\BibleQuote{我从天降下，不是为执行我的旨意，而是为执行派遣我来的父的旨意}；这里他称自己肉身的意志为\Quote{我的旨意}，因为他的肉身也是他自己的。因为正如他至圣无玷、有灵的肉身并未因被神化而毁灭，而是继续存在于其自身的状态和本性内（\OriginalTerm{以其界说和理性}{\GreekText{ὄρῳ} \GreekText{τε} \GreekText{καὶ} \GreekText{λόγῳ}}），同样地，他的人性意志虽被神化，却未被压制，反而得以保存，按照神学家格里高利的说法：\Quote{他的意志（即救主的意志）并不与天主相反，而是完全被神化。}\par

我们在同一我们的主耶稣基督、我们真实的天主内，光荣两个本性的行动，不可分、不变、不混乱、不离散，即一个天主性的行动和一个人的行动，按照神圣的宣讲者良，他极其明确地断言如下：\Quote{因为每一个\InnerQuote{形态}（\OriginalTerm{形态}{\GreekText{μορφὴ}}）与另一个共融时，做其专有之事，即圣言行圣言的事，肉身行肉身的事。}\par

因为我们不允许在天主和受造物中只有一个本性的行动，正如我们不将受造物提升到天主性本质，也不将天主性的荣耀降低到适于受造物的地位。\par

我们承认奇迹和苦难同属一个位格，但由于祂所是并存在于其中的这一性或那一性，正如济利禄精妙所言。因此，保持不混淆与不可分，我们简明地作此整个宣认，相信我们的主耶稣基督是圣三中的一位，并在降生成人后是我们的真天主。我们说，祂的两个性体在祂一个位格中彰显出来，在祂\OriginalTerm{整个救恩经济的交往}{\GreekText{δἰ} \GreekText{ὅλης} \GreekText{αὐτοῦ} \GreekText{τῆς} \GreekText{οἰκονομκῆς} \GreekText{ἀναστροφῆς}}中，祂既施行奇迹，也忍受苦难，并非仅仅是外表，而是真实地，这是由于在同一位置中必须被承认的性体差异；因为虽结合一起，每一性体却愿意并实行其自身之事，且是不可分、不混淆的。因此，我们宣认两个意志和两个运作，在祂内极其适宜地共同协作，以拯救人类。\par

因此，我们极其勤勉审慎地制定了这些事项，兹规定：任何人不得提出、书写、编撰、思想或教导不同的信仰。凡胆敢编撰不同的信仰，或向那些愿意从外邦人、犹太人或任何异端皈依真理认识的人，提出、教导或交付任何不同的信经；或引入新的说法或言辞发明以颠覆我们现在所决定之事者，所有这些，若是主教或圣职人员，应予以撤职：主教免去主教职，圣职人员免去圣职；若是隐修士或平信徒，则应处以绝罚。\par

% structure: p0078 source=h2 target=section reason=relative-heading-rank; base=h2

\section{\WorkTitle{向皇帝的颂词}}

\EditorialNote{这篇讲话以许多对皇帝的赞扬开始，特别是对他为真信仰的热忱。}\par

但由于仇敌撒殚不让人安歇，他兴起了基督的仆役来反对基督，如同武装并携带武器等。\par

\EditorialNote{随后列举了各种异端，以及前五次大公会议如何谴责他们。}\par

事态如此，您这位基督所爱的威严，有必要召集这一至圣而庞大的集会。\par

此后，受圣神感动，全体同意并一致赞同，并认可我们至荣福、至崇高的教宗阿加托的教义信函（他呈递给陛下），同时赞同在他主持下的一百二十五位教父的神圣会议的提议，我们教导：圣三中的一位，我们的主耶稣基督，降生成人，并应在两个完美的性体内，无分裂、无混淆地受到颂扬。因为作为圣言，祂与天主圣父同体、永恒；但作为由无玷童贞玛利亚、天主之母所取的肉身，祂是完美的人，与我们同体，并在时间中受造。因此，我们宣告祂在神性上是完美的，同时同一个祂在人性上也是完美的，按照教父的古老传统和加采东的神圣定义。\par

正如我们承认两个性体，同样我们也承认两个自然的意志和两个自然的运作。因为我们不敢说，在基督内，于祂的降生成人中，任一性体没有意志和运作：恐怕除去那些性体的特性时，我们也除去了它们所归属的性体。因为我们既不否认祂人性的自然意志，也不否认其自然的运作：恐怕我们也否定了为我们救恩而行的经济的主要部分，并将苦难归于神性。这正是那些最近引入可憎的新奇说法（即祂内只有一个意志和一个运作）者所企图做的，他们重振了亚略、阿波利纳里、欧提基和塞维鲁的恶意。因为若我们说主的人性没有意志和运作，我们如何能安全地肯定完美的人性？因为构成人性完整性的，无非是本性的意志，通过此，自由意志的力量在我们内被标明；对于本性的运作也是如此。因为若祂作为人毫无受苦和行动，我们怎能称祂在人性上是完美的？因为正如两个性体的结合为我们保存了一个不混淆、不分裂的位格，同样，这一个位格在两个性体中显现，将各性体所属之事作为自己的事来表明。\par

因此，我们宣告祂内有两个自然的意志和两个自然的运作，共同而无分裂地进行。但我们从教会中驱逐，并理当处以绝罚一切多余的革新及其发明者：即法兰的狄奥多、塞尔吉、保禄、彼洛和伯多禄（他们曾任君士坦丁堡宗主教），此外还有亚历山大里亚的司祭之长基利，以及与他们同流的罗马的\OriginalTerm{治理者}{\GreekText{πρόεδρον}}鸿诺里，正如他在这事上追随了他们。此外，我们理当绝罚并撤职安提约基雅的主教玛加略及其弟子斯德望（或更确切地说，应是其师傅），他试图为其前辈的不虔敬辩护，并且扰乱了整个世界，通过其有害的信函和欺诈的制度，在各地蹂躏了众多人。同样，那位老迈的颇利克罗尼，怀着婴儿般的智力，他曾许诺要复活死人，而当死人未复活时，他受到嘲笑；以及所有曾经教导、正在教导或胆敢教导在降生成人的基督内只有一个意志和一个运作的人。……但宗徒之长与我们一同争战：因为我们有他的效仿者和其宗座的继承人在我们一方，他在其信函中也陈明了\OriginalTerm{神学}{\GreekText{θεολογίας}}的奥秘。因为古老的罗马城向你们递交了一份神圣性格的宣认，从日落之处升起了一份教义图表，使黑暗显明，伯多禄通过阿加托说话，而您，独裁的君王，按照神圣谕令，与您宝座的全能共享者一同裁判。\par

但，仁爱而喜正义的主，求祢回报那位将祢大权赐予祢的人；并请赐下祢御笔签署的谕旨，作为我们所定之事的印证，并以惯常的敬虔谕令和宪章予以确认，使无人能反对已行之事，或再提出任何新问题。因为请放心，温和的陛下，我们并未篡改大公会议及受认可的教父们所定之事，而是予以确认。如今我们同心同声呼喊：\Quote{天主，拯救君王！等等，等等。}\par

\EditorialNote{随后是对皇帝的诸多颂词及为其祈祷的请求。}\par

% structure: p0088 source=h2 target=section reason=relative-heading-rank; base=h2

\section{大公会议致圣阿加托的信}

神圣第六届大公会议致最蒙福、最圣洁的旧罗马教宗阿加托的信函副本。\par

神圣而大公的会议，因天主的恩宠及最虔诚忠信的伟大君士坦丁皇帝的敬虔许可，聚集于这座由天主保护的皇城君士坦丁堡——新罗马，在称为特鲁卢斯的皇宫\OriginalTerm{秘所}{Secretum}（\OriginalTerm{神圣}{\GreekText{θείου}}，\OriginalTerm{（拉丁文）}{sacri}）内，向旧罗马最圣洁、最蒙福的教宗阿加托致意：愿主赐健康。\par

严重的疾病需要更强大的治疗，最蒙福的神父，如您所知；因此，基督我们的真天主，万物的创造者与治理者，赐下一位智慧的医生，即您那受天主尊崇的圣德，以正统信仰的疗法强力驱除异端瘟疫的传染，并给予教会肢体健康的活力。因此，如同对普世教会首席宗座的监督者，我们将应行之事交托于您，因为您乐于以信仰的坚固磐石为立足之地，我们从您以父性真福致最虔诚皇帝的信函中，读到您真实无误的信仰宣认，因而知道这事；我们承认这信函是 \OriginalTerm{神圣地写就}{perscriptas}，如同由宗徒之长所书写；并通过它，我们驱逐了最近兴起的、充满谬误的异端派别，我们受那位神圣统治、手握最温和权杖的君士坦丁所激励而制定法令。藉其帮助，我们推翻了不敬的谬误，仿佛围困了异端者的邪恶教义。然后，我们撕碎他们可憎异端的基础，以属灵与父性的武器攻击他们，混乱他们的口舌，使他们无法彼此协调，推翻了这些最不敬异端追随者所建造的高塔；我们以绝罚处决他们，视为信仰上的堕落者与罪人，在清晨于天主会幕的营外（按达味的话来说），依照您信函中已做出的判决宣告他们：这些名字是：斐兰的主教德奥多禄、塞尔吉乌斯、何诺里乌斯、济利禄、保禄、彼鲁斯与伯多禄。此外，除他们之外，我们也公正地将那些生活于他们所接受的——或更准确地说，这些受天主憎恨之人，即阿波利纳瑞斯、塞维鲁斯与德默斯提乌斯的——不敬之中的人，附以异端者的绝罚，即：安提约基亚大城的玛加略主教（我们也因他对正统信仰的不悔改及顽固执拗，理应剥去他的牧者圣职），以及其疯狂的弟子与不敬的导师斯德望，还有在他异端教义中根深蒂固的颇里略尼禄（恰如其名）；最后，凡是不悔改教导或曾教导、或如今持有或曾持有类似教义的人，均在此列。\par

迄今为止，悲伤、忧愁与多泪是我们分内之事。因为我们不能因邻人的跌倒而欢笑，也不因他们恣意的狂妄而欢呼；我们未曾因此事而自高，以免更加沉重地跌倒；可敬且神圣之首，我们并非如此受教——我们以宇宙的主基督为至仁爱、至爱人者；因祂劝勉我们要在司祭职上，尽可能成为祂的效法者，合乎善人的身份，并效法祂牧者与和好治理的榜样。而且，最温和的皇帝与我们自己，也曾以多种方式劝诫他们悔改，并以极大的虔敬与谨慎处理了全部事务。我们如此行，并非为利益，也非出于仇恨，您可从每次会议中的记录（附函呈奉于您的真福）轻易看出；您也将从您的圣德代表：天主喜爱的司铎德奥多禄与乔治，最虔诚的执事若望，最可敬的副执事君士坦丁（他们都是您的神子，也是我们亲爱的弟兄）那里得知。同样，您也会从您神圣主教会议所派遣的圣主教们那里听到同样的事，他们按照您的纪律，正确正直地与我们一同在信仰的第一章中制定了决议。\par

如此，蒙圣神光照，受您教义指引，我们抛出了不敬的邪恶教义，开辟了正统信仰的正直道路，并因我们最虔诚温和的皇帝君士坦丁的智慧与力量，在各方面得到鼓励与协助。然后，我们中的一员——这首都君士坦丁堡最神圣的\OriginalTerm{主教}{praesul}——首先赞同您致最虔诚皇帝的正统信函，认为它完全符合经认可的教父及受天主教导的教父、以及神圣五次大公会议的教导；我们全体，因基督我们天主的助佑，轻易完成了我们所追求的事。因为天主既是推动者，天主也加冕了我们的会议。\par

是故，圣神的恩宠照耀我们，借着你们恳切的祈祷显示其德能，为拔除一切蒺藜和不结好果子的树，并吩咐将它们用火焚烧。我们众人皆同心、同口、同手，赖赋予生命的圣神之助，颁布了一项纯洁无误、确定无误的定断；并非如经上所记\Quote{移去古时的地标}（愿天主禁止！），而是坚守圣洁而受认可的教父们的见证与权威，并明定：如同基督我们的真天主，是由两个性体（即神性与人性）所组成，并存在于其中，我们宣讲他、荣耀他，是不可分、不变、不混、不分离的；同样，我们也承认他有两个本然的活动，是不分、不争、不混、不分离的，正如我们的会议定断所宣告的。这些谕令得到了我们效法天主的皇帝陛下的同意，并由他亲手签署。且如所述，我们弃绝并谴责那最不敬、无实的异端，它主张在我们降生成人的真天主基督内只有一个意志和一个活动；借此我们严厉压制了那些混淆与分裂之人，扑灭了其他异端燃起的风暴，而与你们一起清楚地展现了正统信仰的明光；并恳请你们的神圣父爱，以尊贵的覆函确认我们的定断；借此我们满怀在基督内的希望，深信他的仁慈恩惠将赐予罗马帝国（托付于我们最仁慈的皇帝治理）安定；并以日复一日的胜利装饰他极温和的仁政；且在他今世赐予我们的恩惠之外，将你们那受天主尊崇的圣德置于他可畏的审判台前，作为一位真诚宣认真信仰、保存其无玷、并忠于托付于他的正统羊群的好牧者。\par

我们及所有与我们同在的人，问候与你们真福同在的基督内众弟兄。\par

% structure: p0096 source=h2 target=section reason=relative-heading-rank; base=h2

\section{皇帝敕令，张贴于大教堂第三门厅，近所谓戟钹处}

因我们的主、导师——天主及救主耶稣基督之名，至为虔诚的皇帝，平和而爱基督的君士坦丁，在耶稣基督内忠于天主的皇帝，致住在由天主护守的皇城内的所有爱基督的子民。\par

\Emph{该文献极长，赫费勒给出了以下摘要，对普通读者而言已足够，读者须记得这是皇帝的敕令，而非大公会议的文件。}\par

\Emph{赫费勒的摘要}（\WorkTitle{会议史}，卷五，第178页）。\par

\Quote{阿波林纳里等的异端，由法兰的德奥多若复兴，并由旧罗马教宗何诺留（他曾自相矛盾）所确认。此外，居鲁士、皮尔若斯、保禄、伯多禄；更近的有玛加略、斯德望及波利克罗尼乌斯，均传播了唯意志论。因此，皇帝召集了此神圣大公会议，并颁布本敕令及信仰明辨，以确认并巩固其定断。（随后是一段详尽的信仰明辨，并引证两个意志和活动的道理。）正如他承认前五次大公会议，同样他也弃绝从西满魔术师起的所有异端者，尤其此新异端的发起人与赞助者：德奥多若与塞尔基乌斯；还有教宗何诺留，\OriginalTerm{他在一切事上同为他们的盟友、同谋并肯定异端者}{\GreekText{τὸν} \GreekText{κατὰ} \GreekText{πάντα} \GreekText{τούτοις} \GreekText{συναιρέτην} \GreekText{καὶ} \GreekText{σύνδρομον} \GreekText{καὶ} \GreekText{βεβαιωτὴν} \GreekText{τῆς} \GreekText{αἱρέσεως}}，以及居鲁士等人，并规定此后任何人不得持守不同信仰，不得妄图教导一个意志和一个活动。人若非在正统信仰内，不能得救。凡不服从皇帝敕令者，如系主教或神职人员，予以撤职；如系官员，则没收财产并夺去束带（\OriginalTerm{腰带}{\GreekText{ζώνη}}）；如系平民，则流放离宫及所有城市。}\par

\SourceRecord{Translated by Henry Percival.；Nicene and Post-Nicene Fathers, Second Series；Vol. 14.；Edited by Philip Schaff and Henry Wace.；Buffalo, NY: Christian Literature Publishing Co.,；1900.；<http://www.newadvent.org/fathers/3813.htm>.}
